\chapter{Tổng quan đề tài}

\section{Bối cảnh}

Dealzy là hệ thống thương mại điện tử tập trung vào mô hình bán voucher giảm giá trực tuyến. Hệ thống đóng vai trò trung gian giữa khách hàng có nhu cầu mua ưu đãi và các đối tác cung cấp dịch vụ như nhà hàng, du lịch, spa, mua sắm hoặc giải trí. Trong phạm vi đồ án, nhóm xây dựng một nền tảng demo có đủ ba vai trò cốt lõi: Khách hàng, Đối tác và Quản trị viên.

Điểm trọng tâm của đề tài không chỉ là hiển thị voucher, mà còn là quản lý vòng đời voucher: Đối tác tạo voucher, Quản trị viên kiểm duyệt, voucher được công bố cho Khách hàng, Khách hàng mua voucher và Đối tác xác thực việc sử dụng voucher. Đến thời điểm báo cáo, các luồng quản lý, kiểm duyệt, mua voucher mô phỏng, phát hành E-Voucher và xác thực voucher đã được hiện thực ở mức có thể chạy demo end-to-end.

\section{Mục tiêu}

\begin{itemize}[leftmargin=1.2cm]
    \item Xây dựng hệ thống thương mại điện tử có cơ sở dữ liệu quan hệ và phân quyền theo vai trò.
    \item Cung cấp giao diện riêng cho khách hàng, đối tác và quản trị viên để giảm chồng chéo nghiệp vụ.
    \item Thiết kế backend theo kiến trúc module để mỗi nhóm vai trò có route, controller và service riêng.
    \item Hiện thực các quy tắc nghiệp vụ quan trọng bằng cả code backend và ràng buộc trong PostgreSQL.
    \item Chuẩn bị bộ sản phẩm bàn giao gồm báo cáo, sơ đồ, ERD, script database, dữ liệu mẫu và mã nguồn ứng dụng khi chốt bản ổn định.
\end{itemize}

\section{Phạm vi hiện tại}

Phạm vi hiện tại của Dealzy bao phủ ba vai trò chính gồm Khách hàng, Đối tác và Quản trị viên. Các chức năng tra cứu voucher, quản lý voucher, kiểm duyệt, mua hàng mô phỏng, phát hành E-Voucher, xác thực mã, khiếu nại, đánh giá, nội dung, nhật ký, chatbot và cập nhật realtime đã có trong mã nguồn. Quy trình Order/Payment sử dụng cổng thanh toán mô phỏng/sandbox, tạo đơn trong transaction, chuyển đơn sang Paid khi có xác nhận hợp lệ và phát hành E-Voucher cho từng đơn vị voucher đã mua.

\begin{table}[H]
\centering
\small
\begin{tabularx}{\textwidth}{|p{3.0cm}|Y|p{3.2cm}|}
\hline
\textbf{Mảng} & \textbf{Nội dung} & \textbf{Trạng thái} \\ \hline
Khách hàng & Đăng ký/đăng nhập, hồ sơ cá nhân, tìm kiếm/lọc voucher, xem chi tiết, giỏ hàng, checkout, trạng thái thanh toán, lịch sử đơn, ví E-Voucher/QR, đánh giá, khiếu nại và chatbot & Đã có \\ \hline
Đối tác & Đăng ký doanh nghiệp, quản lý hồ sơ và chi nhánh, dashboard, tạo/cập nhật/gửi duyệt/tạm ngưng voucher, xem báo cáo, tra cứu và xác thực mã voucher bằng mã/QR, nhận cập nhật realtime & Đã có \\ \hline
Quản trị viên & Duyệt/từ chối đối tác, quản lý người dùng, duyệt/từ chối/ẩn voucher, dashboard, quản lý đơn hàng, xác nhận thanh toán thủ công và hoàn tiền mô phỏng & Đã có \\ \hline
Mua hàng và thanh toán & Checkout, kiểm tra tồn kho, tạo Orders/Order\_Items, thanh toán mô phỏng qua VietQR/MoMo/PayPal và phát hành E-Voucher sau khi đơn được xác nhận Paid & Đã có \\ \hline
Dữ liệu & PostgreSQL, ERD, script tạo bảng, seed data, ràng buộc khóa ngoại/UNIQUE/CHECK, trigger kiểm tra đơn hàng, xác thực sử dụng voucher, kiểm tra điều kiện đánh giá & Đã có \\ \hline
Bảo mật và vận hành & JWT, bcrypt, middleware phân quyền theo vai trò, CORS, Helmet, rate limit, health check API/database, cấu hình deploy Render/Vercel & Đã có\\ \hline
\end{tabularx}
\caption{Tổng quan phạm vi hiện thực}
\end{table}
\section{Công nghệ sử dụng}

\begin{itemize}[leftmargin=1.2cm]
    \item Frontend: React, Vite, React Router, Framer Motion, Lucide React, Recharts.
    \item Backend: Node.js, Express, JWT, bcryptjs, pg, Nodemailer.
    \item Database: PostgreSQL.
    \item Tài liệu: LaTeX, Mermaid, Markdown.
    \item Quản lý mã nguồn: Git, GitHub; nhánh tổng hợp hiện tại là \texttt{partner}.
\end{itemize}
